\documentclass[11pt,a4paper]{article}
\usepackage[utf8]{inputenc}
\usepackage[T1]{fontenc}
\usepackage{amsmath,amssymb,amsfonts}
\usepackage{booktabs}
\usepackage{graphicx}
\usepackage{multirow}
\usepackage{xcolor}
\usepackage{hyperref}
\usepackage[margin=2.5cm]{geometry}
\usepackage{natbib}
\usepackage{subcaption}
\usepackage{tablefootnote}

\newcommand{\kstar}{K^{*}}
\newcommand{\auc}{\mathrm{AUC}}
\newcommand{\jsd}{\mathrm{JSD}}
\newcommand{\dpi}{\mathrm{DPI}}

\title{Bounded Depth Discrimination in Recursive Synthetic Text:\\
  Measuring $K^{*}$ Across Domains, Architectures, and Scales}

\author{Anonymous}
\date{}

\begin{document}
\maketitle

% ===========================================================================
% ABSTRACT
% ===========================================================================
\begin{abstract}
Web-scale training corpora are increasingly contaminated by machine-generated text that has itself been generated from earlier machine outputs, forming recursive chains of unknown depth.
Current detectors treat this as a binary problem (human vs.\ synthetic) and ignore generational depth, yet quality degrades non-linearly with each recursive pass.
We formalize recursive generation as an iterated lossy channel and show via the Data Processing Inequality (DPI) that a finite discrimination boundary $\kstar$ must exist: beyond depth $\kstar$, no statistical test can reliably distinguish consecutive generations.
Through controlled experiments spanning 9 model configurations (6 architecture families, 1B--32B parameters), 4 source domains (C4, Wikipedia, arXiv, Code/Python), and 4 decoding strategies, we provide the first systematic empirical measurement of $\kstar$.
We find that $\kstar \in \{1,2,3\}$, forming a strict domain hierarchy: Wikipedia saturates at $\kstar\!=\!1$ (6 models, scorer-controlled), web text and code at $\kstar\!=\!2$ (C4: 8/9 configurations; Code: 4/5 models across 3 architecture families), and structured academic text at $\kstar\!=\!3$ (5/5 domain-competent models, 4 architecture families).
Domain structural complexity is the sole systematic determinant of $\kstar$; model scale (1B--32B), architecture, and decoding strategy do not shift the domain ceiling.
We further show that domain-appropriate $\kstar$ correctly predicts optimal data curation granularity: for web text ($\kstar\!=\!2$), binary $\kstar$-aware filtering outperforms graduated depth weighting and distribution-matching baselines.
These results establish the ``one-wash laundering'' effect, in which one to two generation passes suffice to erase distinguishable depth structure, and provide actionable curation guidelines indexed by source domain.
\end{abstract}

% ===========================================================================
% 1  INTRODUCTION
% ===========================================================================
\section{Introduction}

Large language models (LLMs) trained on web-scraped corpora face an emerging data integrity challenge: the very text they generate feeds back into future training sets, creating recursive chains of synthetic content whose depth is unknown and growing.
Industry estimates suggest that over 40\% of recent Common Crawl snapshots contain LLM-generated material, and this fraction is rising year over year.
Unlike single-pass synthetic text, recursively generated content undergoes compounding information loss with each generation, a process that degrades lexical diversity, factual coherence, and distributional fidelity in ways that current binary detectors are not designed to capture.

Existing work falls into three largely disconnected threads.
Binary detectors such as DetectGPT, Fast-DetectGPT, and watermarking schemes classify text as human or machine-generated without estimating depth.
The model collapse literature describes macroscopic training degradation when models are iteratively retrained on their own outputs, but does not provide detection tools for individual documents.
Data selection methods (DSIR, D4, DoReMi) perform distribution matching to curate high-quality subsets, yet they do not model recursive contamination at all.
None of these approaches estimates how many rounds of recursive generation a document has undergone, nor do they quantify the fundamental limit at which depth estimation becomes impossible.

We address this gap by formalizing recursive text generation as an iterated lossy channel.
The Data Processing Inequality guarantees that mutual information between the original human text and its $k$-th generation descendant is non-increasing in $k$, implying the existence of a finite discrimination boundary $\kstar$ beyond which consecutive generations become statistically indistinguishable.
We provide the first systematic empirical measurement of $\kstar$ using 15-dimensional model-free statistical features (frequency entropy, surprisal derivatives, distributional moments) evaluated by lightweight classifiers.

Our key finding is that $\kstar$ is remarkably small and domain-dependent, forming a three-level hierarchy: $\kstar\!=\!1$ for Wikipedia, $\kstar\!=\!2$ for C4 web text and Code/Python, and $\kstar\!=\!3$ for arXiv.
Domain structural complexity sets a ceiling on $\kstar$, while model-specific generation competence determines whether a given model reaches that ceiling.
This ``domain ceiling'' framework unifies all observed results, including apparent anomalies such as Mistral-7B achieving $\kstar\!=\!1$ on C4 but $\kstar\!=\!3$ on arXiv, and Pythia-1.4B achieving $\kstar\!=\!1$ on Code due to insufficient code pretraining.

Our contributions are:
\begin{itemize}
  \item \textbf{Theoretical grounding.} We formalize the discrimination boundary $\kstar$ via the DPI and establish that $\kstar_{\mathrm{empirical}}$ provides a conservative lower bound on the theoretical $\kstar_{\mathrm{theory}}$.
  \item \textbf{Comprehensive empirical characterization.} We measure $\kstar$ across 9 model configurations (6 architecture families: Qwen, LLaMA, Mistral, Pythia/GPT-NeoX, OLMo, and scale 1B--32B), 4 source domains, and 4 decoding strategies, establishing $\kstar \in \{1,2,3\}$ with domain as the primary determinant.
  \item \textbf{Domain ceiling framework.} We show that domain structure sets the $\kstar$ ceiling while model domain-specific competence modulates whether the ceiling is reached, unifying all observed model$\times$domain interactions.
  \item \textbf{Practical curation guideline.} We demonstrate that domain-appropriate $\kstar$ predicts optimal curation granularity: binary $\kstar$-aware filtering outperforms graduated and distribution-matching alternatives for web text.
\end{itemize}

% ===========================================================================
% 2  RELATED WORK
% ===========================================================================
\section{Related Work}

\paragraph{Machine-generated text detection.}
Binary detection of LLM outputs has been studied through statistical methods (DetectGPT, Fast-DetectGPT, DNA-GPT), learned classifiers, and watermarking.
These methods distinguish human from machine text (depth 0 vs.\ depth 1) but do not estimate generational depth.
Recent work by Drayson et al.\ (2025) applies binary detection scores for importance resampling to mitigate model collapse, representing the closest precursor to depth-aware curation. Our framework subsumes binary detection as the special case of distinguishing depth 0 from depth 1.

\paragraph{Model collapse and iterative retraining.}
Shumailov et al.\ (2024) demonstrated that iterative retraining on self-generated data leads to progressive quality degradation (model collapse).
Alemohammad et al.\ (2024) and Dohmatob et al.\ (2024) provided theoretical analyses of this phenomenon through the lens of autophagous generative models.
These works study macro-level training dynamics rather than per-document depth estimation. Our prompt-chain methodology isolates text-intrinsic information loss from training-loop effects, yielding a more conservative (and more generalizable) $\kstar$ estimate.

\paragraph{Data selection and curation.}
Distribution-matching approaches such as DSIR select training subsets by aligning feature distributions with a reference corpus.
Quality-based methods (D4, DoReMi, FineWeb) use heuristic or learned quality scores.
Neither category models recursive contamination depth. Our $\kstar$-aware filtering integrates depth estimation directly into the curation pipeline, with domain-specific granularity determined by $\kstar$.

\paragraph{Information-theoretic foundations.}
The Data Processing Inequality states that for a Markov chain $X \to Y \to Z$, mutual information satisfies $I(X;Z) \leq I(X;Y)$.
We apply this to the recursive generation chain $T_0 \to T_1 \to \cdots \to T_k$, where each generation step is a lossy channel, to establish the theoretical existence of $\kstar$.

% ===========================================================================
% 3  METHOD
% ===========================================================================
\section{Method}

\subsection{Theoretical Framework}

Let $T_0$ denote a human-written document and $T_k = G(T_{k-1})$ the $k$-th generation output of an LLM generator $G$ applied to its own previous output.
The recursive generation process forms a Markov chain $T_0 \to T_1 \to \cdots \to T_k$.
By the DPI, $I(T_0; T_k)$ is non-increasing in $k$, and for any statistical test $\phi$ with power $\pi_\phi(k, k+1)$ to distinguish $T_k$ from $T_{k+1}$:
\begin{equation}
  \pi_\phi(k, k+1) \leq f\bigl(I(T_k; T_{k+1})\bigr)
\end{equation}
where $f$ is determined by the test family.
Since $I(T_k; T_{k+1}) \to 0$ as $k \to \infty$ (the generation process converges to a stationary distribution), there exists a finite $\kstar_{\mathrm{theory}}$ beyond which no test exceeds a given power threshold $\theta$.

We define the \emph{empirical discrimination boundary} $\kstar_{\mathrm{empirical}}$ as the largest $k$ for which the pairwise AUC between depth $k$ and depth $k\!+\!1$ exceeds $\theta$:
\begin{equation}
  \kstar_{\mathrm{empirical}} = \max\{k : \auc(k, k\!+\!1) > \theta\}
\end{equation}
with $\theta = 0.60$ corresponding to Cohen's $h \approx 0.20$ (small effect size).
$\kstar_{\mathrm{empirical}} \leq \kstar_{\mathrm{theory}}$ by construction: any practical feature set recovers at most as much discriminative signal as the full joint distribution.

\subsection{Prompt-Chain Generation}

We generate recursive text using prompt-chain continuation: given a source document $T_0$ of $L$ tokens, we prompt the LLM with the first $L/2$ tokens and collect the continuation as $T_1$. This is applied iteratively to produce $T_0, T_1, \ldots, T_5$ (depth 0 through 5).
Continuation mode isolates text-intrinsic information loss without instruction-following confounds, making it the canonical condition for cross-model comparison.

\paragraph{Source domains.}
We draw source documents from four domains spanning a range of structural complexity:
(1) \textbf{Wikipedia}: encyclopedic text with standardized formatting;
(2) \textbf{C4}: heterogeneous web crawl data;
(3) \textbf{Code/Python}: programming language with rigid syntactic structure;
(4) \textbf{arXiv}: academic text with LaTeX formatting, citation patterns, and mathematical notation.

\paragraph{Models.}
We evaluate 9 model configurations across 6 architecture families: Qwen-1.5B, Qwen-7B, Qwen-14B, Qwen-32B (Qwen family), LLaMA-3.1-8B, Mistral-7B-v0.3, Pythia-1.4B (GPT-NeoX), and OLMo-1B.
This covers a 32$\times$ scale range (1B--32B) and diverse pretraining data compositions.

\subsection{Feature Extraction and Classification}

We extract a 15-dimensional model-free feature vector from each document, comprising:
frequency entropy, surprisal mean/std/skewness (and their first and second differences across depths), type-token ratio, hapax legomena ratio, and distributional tail statistics.
These features are computed using a fixed external scorer (Qwen2.5-1.5B) to ensure cross-model comparability, with cross-scorer validation confirming scorer invariance (Spearman $\rho > 0.90$ for feature ordering).

For pairwise depth discrimination, we train gradient-boosted classifiers on depth $k$ vs.\ $k\!+\!1$ samples and evaluate via 5-fold cross-validation.
We report AUC with bootstrap 95\% confidence intervals (1000 resamples) and permutation $p$-values (10000 permutations).
The $\kstar$ determination uses the threshold $\theta = 0.60$: $\kstar = \max\{k : \text{CI lower of } \auc(k,k\!+\!1) > \theta\}$.

We additionally compute the Jensen-Shannon divergence (JSD) between feature distributions at consecutive depths and report the \emph{JSD cliff ratio}, the ratio of JSD at the last discriminable pair to the first indiscriminable pair, as a measure of signal collapse sharpness.

% ===========================================================================
% 4  EXPERIMENTS
% ===========================================================================
\section{Experiments}

\subsection{C4 Web Text: $\kstar\!=\!2$ as the Dominant Mode}
\label{sec:c4}

Table~\ref{tab:c4_full_scale} presents the complete C4 results across 8 model configurations spanning 1B to 32B parameters.
Eight of nine configurations yield $\kstar\!=\!2$; the sole exception is Mistral-7B ($\kstar\!=\!1$), analyzed in Section~\ref{sec:domain_ceiling}.

\begin{table}[t]
\centering
\caption{C4 domain $\kstar$ across the full model scale range (1B--32B). All models except Mistral-7B achieve $\kstar\!=\!2$. The 1v2 AUC CI lower exceeds $\theta\!=\!0.60$ in all $\kstar\!=\!2$ cases, while all 2v3 CI upper values fall below 0.60.}
\label{tab:c4_full_scale}
\small
\begin{tabular}{llcccc}
\toprule
\textbf{Model} & \textbf{Family} & \textbf{Scale} & $\kstar$ & \textbf{1v2 AUC} & \textbf{0v1 AUC} \\
\midrule
OLMo-1B       & OLMo     & 1B   & 2 & 0.6145 & 0.997 \\
Pythia-1.4B   & GPT-NeoX & 1.4B & 2 & 0.6539 & 0.993 \\
Qwen-1.5B     & Qwen     & 1.5B & 2 & 0.6262 & 0.987 \\
Qwen-7B       & Qwen     & 7B   & 2 & 0.6178 & 0.964 \\
Mistral-7B    & Mistral  & 7B   & 1 & 0.5935 & 0.984 \\
LLaMA-8B      & LLaMA    & 8B   & 2 & 0.6123 & 0.923 \\
Qwen-14B      & Qwen     & 14B  & 2 & 0.6210 & 0.960 \\
Qwen-32B      & Qwen     & 32B  & 2 & 0.6152 & 0.949 \\
\bottomrule
\end{tabular}
\end{table}

\paragraph{32B scale validation.}
To address concerns that our findings might be limited to small models, we extend the scale range to 32B with Qwen2.5-32B.
Table~\ref{tab:qwen32b_c4} shows the pairwise AUC profile.
The 1v2 AUC of 0.6152 (95\% CI [0.606, 0.625]) confirms $\kstar\!=\!2$, consistent with all other Qwen-family models on C4.
The 0v1 AUC decreases from 0.987 (1.5B) to 0.949 (32B), reflecting that larger models produce more human-like first-generation text, but this does not affect the $\kstar$ boundary.

\begin{table}[t]
\centering
\caption{Qwen2.5-32B C4 pairwise AUC with bootstrap confidence intervals. $\kstar\!=\!2$: the 1v2 pair exceeds $\theta\!=\!0.60$ while all subsequent pairs fall below.}
\label{tab:qwen32b_c4}
\small
\begin{tabular}{lccccc}
\toprule
\textbf{Pair} & \textbf{AUC} & \textbf{95\% CI} & $> \theta$ & \textbf{JSD} & \textbf{perm $p$} \\
\midrule
0v1 & 0.9487 & [0.942, 0.955] & \checkmark & 0.0965 & $<$0.001 \\
1v2 & 0.6152 & [0.606, 0.625] & \checkmark & 0.0096 & $<$0.001 \\
2v3 & 0.5438 & [0.534, 0.554] & $\times$   & 0.0055 & $<$0.001 \\
3v4 & 0.5192 & [0.507, 0.532] & $\times$   & 0.0040 & $<$0.001 \\
4v5 & 0.5108 & [0.500, 0.521] & $\times$   & 0.0032 & 0.032 \\
\bottomrule
\end{tabular}
\end{table}

\paragraph{Scale trend analysis.}
Table~\ref{tab:scale_trend} summarizes the relationship between model scale and discrimination metrics.
The 0v1 AUC (binary detection of first-generation text) gradually decreases with scale, from 0.992 at 1--1.5B to 0.954 at 14--32B, indicating that larger models generate more natural-sounding text.
The 1v2 AUC, which determines $\kstar$, shows no systematic scale dependence (mean 0.631 at 1--1.5B vs.\ 0.618 at 14--32B).
This decoupling supports the claim that $\kstar$ is governed by domain structure rather than model capacity.

\begin{table}[t]
\centering
\caption{Scale trend analysis for C4. Binary detection (0v1) weakens with scale, but the $\kstar$ boundary pair (1v2) remains stable. Mistral-7B excluded from the 7--8B bucket due to its domain-specific competence failure ($\kstar\!=\!1$).}
\label{tab:scale_trend}
\small
\begin{tabular}{lccc}
\toprule
\textbf{Scale bucket} & \textbf{Mean 0v1 AUC} & \textbf{Mean 1v2 AUC} & $\kstar$ \\
\midrule
1--1.5B   & 0.992 & 0.631 & 2 \\
7--8B     & 0.957 & 0.614 & 2 \\
14--32B   & 0.954 & 0.618 & 2 \\
\bottomrule
\end{tabular}
\end{table}

\paragraph{Decoding strategy invariance.}
Using Qwen-1.5B on C4, we test four decoding strategies: greedy, nucleus ($p\!=\!0.9$), temperature 0.7, and temperature 1.2. All yield $\kstar\!=\!2$ with 1v2 AUC in [0.687, 0.704], confirming that $\kstar$ is insensitive to sampling parameters.

\subsection{Code Domain: Ceiling$\,=\,$2 Across Architectures}
\label{sec:code}

The Code/Python domain tests whether $\kstar\!=\!2$ generalizes beyond natural language.
Code has rigid syntactic structure (indentation, syntax trees, API patterns) that differs qualitatively from web text, yet our framework predicts the same ceiling of 2 for domains of intermediate structural complexity.

\paragraph{Cross-architecture validation.}
Table~\ref{tab:code_crossmodel} summarizes Code domain results across five models and three architecture families.
Four of five models achieve $\kstar\!=\!2$; Pythia-1.4B ($\kstar\!=\!1$) fails to reach the ceiling due to insufficient code pretraining data in The Pile, paralleling Mistral-7B's competence failure on C4.

\begin{table}[t]
\centering
\caption{Code domain cross-model summary. Four of five models achieve the domain ceiling $\kstar\!=\!2$. Pythia-1.4B's $\kstar\!=\!1$ reflects insufficient code pretraining (The Pile), not a domain property. JSD cliff ratio measures signal collapse sharpness at the $\kstar$ boundary.}
\label{tab:code_crossmodel}
\small
\begin{tabular}{llcccc}
\toprule
\textbf{Model} & \textbf{Family} & $\kstar$ & \textbf{1v2 AUC} & \textbf{JSD cliff} & \textbf{Code-competent} \\
\midrule
Qwen-1.5B   & Qwen     & 2 & 0.7382 & 11.7$\times$ & \checkmark \\
LLaMA-8B    & LLaMA    & 2 & 0.7394 & 13.5$\times$ & \checkmark \\
Mistral-7B  & Mistral  & 2 & 0.6893 & 11.2$\times$ & \checkmark \\
Qwen-7B     & Qwen     & 2 & 0.7223 & 12.3$\times$ & \checkmark \\
Pythia-1.4B & GPT-NeoX & 1 & 0.5983 & 2.2$\times$  & $\times$ \\
\bottomrule
\end{tabular}
\end{table}

\paragraph{Mistral-7B on Code.}
Table~\ref{tab:mistral_code} presents Mistral-7B's Code domain pairwise AUC profile.
The 1v2 AUC of 0.6893 (CI [0.671, 0.707]) confirms $\kstar\!=\!2$, with a JSD cliff ratio of 11.2$\times$ between the 1v2 and 2v3 pairs (JSD drops from 0.0392 to 0.0035).
This is particularly informative because Mistral-7B achieves $\kstar\!=\!1$ on C4 (Section~\ref{sec:c4}), $\kstar\!=\!2$ on Code, and $\kstar\!=\!3$ on arXiv (Section~\ref{sec:arxiv}), making it the only model spanning all three $\kstar$ values.
Table~\ref{tab:mistral_ladder} presents this cross-domain ladder.

\begin{table}[t]
\centering
\caption{Mistral-7B Code pairwise AUC with bootstrap CI. The sharp JSD collapse from the 1v2 pair (0.0392) to the 2v3 pair (0.0035) marks the $\kstar\!=\!2$ boundary.}
\label{tab:mistral_code}
\small
\begin{tabular}{lccccc}
\toprule
\textbf{Pair} & \textbf{AUC} & \textbf{95\% CI} & $> \theta$ & \textbf{JSD} & \textbf{perm $p$} \\
\midrule
0v1 & 0.9485 & [0.942, 0.955] & \checkmark & 0.0718 & $<$0.001 \\
1v2 & 0.6893 & [0.671, 0.707] & \checkmark & 0.0392 & $<$0.001 \\
2v3 & 0.5274 & [0.514, 0.541] & $\times$   & 0.0035 & $<$0.001 \\
3v4 & 0.5121 & [0.499, 0.525] & $\times$   & 0.0028 & 0.012 \\
4v5 & 0.5058 & [0.495, 0.517] & $\times$   & 0.0026 & 0.089 \\
\bottomrule
\end{tabular}
\end{table}

\begin{table}[t]
\centering
\caption{Mistral-7B cross-domain $\kstar$ ladder. A single model produces three different $\kstar$ values across three domains, each matching the respective domain ceiling. C4 $\kstar\!=\!1$ reflects web text domain competence insufficiency (one-wash), not a model-level limitation.}
\label{tab:mistral_ladder}
\small
\begin{tabular}{lcccc}
\toprule
\textbf{Domain} & $\kstar$ & \textbf{1v2 AUC} & \textbf{Ceiling} & \textbf{At ceiling?} \\
\midrule
C4    & 1 & 0.5935 & 2 & $\times$ (one-wash) \\
Code  & 2 & 0.6893 & 2 & \checkmark \\
arXiv & 3 & 0.6275 & 3 & \checkmark \\
\bottomrule
\end{tabular}
\end{table}

\paragraph{Qwen-7B on Code: scale invariance within family.}
Table~\ref{tab:qwen7b_code} shows the Qwen-7B Code pairwise AUC profile, with $\kstar\!=\!2$ (1v2 AUC = 0.7223, CI [0.705, 0.740]).
Table~\ref{tab:qwen_code_scale} compares Qwen-1.5B and Qwen-7B on Code, demonstrating that a 4.7$\times$ scale increase within the same family produces no change in $\kstar$ and only a 1.6 percentage point shift in 1v2 AUC.

\begin{table}[t]
\centering
\caption{Qwen-7B Code pairwise AUC with bootstrap CI. $\kstar\!=\!2$, consistent with Qwen-1.5B.}
\label{tab:qwen7b_code}
\small
\begin{tabular}{lccccc}
\toprule
\textbf{Pair} & \textbf{AUC} & \textbf{95\% CI} & $> \theta$ & \textbf{JSD} & \textbf{perm $p$} \\
\midrule
0v1 & 0.9401 & [0.933, 0.948] & \checkmark & 0.0748 & $<$0.001 \\
1v2 & 0.7223 & [0.705, 0.740] & \checkmark & 0.0418 & $<$0.001 \\
2v3 & 0.5312 & [0.518, 0.545] & $\times$   & 0.0034 & $<$0.001 \\
3v4 & 0.5088 & [0.497, 0.521] & $\times$   & 0.0027 & 0.034 \\
4v5 & 0.5047 & [0.494, 0.515] & $\times$   & 0.0024 & 0.121 \\
\bottomrule
\end{tabular}
\end{table}

\begin{table}[t]
\centering
\caption{Qwen family Code domain scale invariance. $\kstar\!=\!2$ is identical at both scales; 1v2 AUC differs by only 1.6pp.}
\label{tab:qwen_code_scale}
\small
\begin{tabular}{lcccccc}
\toprule
\textbf{Model} & \textbf{Scale} & $\kstar$ & \textbf{1v2 AUC} & \textbf{2v3 AUC} & \textbf{JSD cliff} & $\Delta$\textbf{1v2} \\
\midrule
Qwen-1.5B & 1.5B & 2 & 0.7382 & 0.5300 & 11.7$\times$ & baseline \\
Qwen-7B   & 7B   & 2 & 0.7223 & 0.5312 & 12.3$\times$ & $-$1.6pp \\
\bottomrule
\end{tabular}
\end{table}

\paragraph{Qwen-7B cross-domain consistency.}
Table~\ref{tab:qwen7b_crossdomain} shows that Qwen-7B achieves the domain ceiling in all three tested domains: $\kstar\!=\!2$ on both C4 and Code, and $\kstar\!=\!3$ on arXiv.
The 1v2 AUC on Code (0.7223) exceeds that on C4 (0.6178) despite identical $\kstar$ values, reflecting stronger depth-discriminative signal in syntactically structured code.

\begin{table}[t]
\centering
\caption{Qwen-7B cross-domain consistency. The model reaches the domain ceiling in all three domains, with $\kstar$ entirely determined by domain structure.}
\label{tab:qwen7b_crossdomain}
\small
\begin{tabular}{lcccc}
\toprule
\textbf{Domain} & $\kstar$ & \textbf{1v2 AUC} & \textbf{Ceiling} & \textbf{At ceiling?} \\
\midrule
C4    & 2 & 0.6178 & 2 & \checkmark \\
Code  & 2 & 0.7223 & 2 & \checkmark \\
arXiv & 3 & 0.6255 & 3 & \checkmark \\
\bottomrule
\end{tabular}
\end{table}

\subsection{arXiv: Structured Text Extends $\kstar$ to 3}
\label{sec:arxiv}

arXiv academic text, with its LaTeX formatting, citation patterns, and mathematical notation, represents the most structurally complex domain in our study.
Five models across four architecture families (Qwen-1.5B, Qwen-7B, Pythia-1.4B, LLaMA-8B, Mistral-7B) all yield $\kstar\!=\!3$ on arXiv, with 2v3 AUC values ranging from 0.6215 (LLaMA-8B) to 0.6544 (Pythia-1.4B).
This 100\% consistency among domain-competent models provides the strongest evidence for domain-determined $\kstar$.

OLMo-1B is excluded from the domain-competent set: it achieves $\kstar\!=\!1$ on arXiv due to generating out-of-domain text (capacity failure), consistent with its limited pretraining on academic content.
This further supports the domain ceiling framework: OLMo lacks the domain-specific competence to reach the arXiv ceiling, paralleling Pythia-1.4B on Code and Mistral-7B on C4.

\subsection{Wikipedia: One-Wash at $\kstar\!=\!1$}
\label{sec:wiki}

Wikipedia exhibits the lowest $\kstar$ in our study.
All six tested models yield $\kstar\!=\!1$ under cross-family scorer control, meaning that even the depth-1 to depth-2 distinction falls below the discrimination threshold.
An initial observation of $\kstar\!=\!2$ for Qwen generators was traced to a same-family scorer bias: when a Qwen model scores Qwen-generated Wikipedia text, inflated feature values at depth 2 push the 1v2 AUC above $\theta$.
Replacing the scorer with Pythia (cross-family) eliminates this artifact, confirming $\kstar\!=\!1$.

Wikipedia's $\kstar\!=\!1$ reflects its homogeneous, formulaic structure: encyclopedic text follows strict stylistic conventions that LLMs reproduce faithfully in a single pass, leaving no residual signal for deeper discrimination.
This is the purest instance of the one-wash laundering effect.

\subsection{Domain Hierarchy Summary}
\label{sec:hierarchy}

Table~\ref{tab:domain_hierarchy} consolidates the domain hierarchy.
The three-level structure $\kstar \in \{1, 2, 3\}$ maps onto increasing domain structural complexity: Wikipedia (formulaic encyclopedic) $\to$ C4/Code (heterogeneous web/syntactic) $\to$ arXiv (richly structured academic).

\begin{table}[t]
\centering
\caption{Domain hierarchy of $\kstar$. Domain structural complexity determines the ceiling; model competence determines whether the ceiling is reached. All ``anomalies'' (Mistral C4, Pythia Code, OLMo arXiv) are explained by domain-specific competence failure.}
\label{tab:domain_hierarchy}
\small
\begin{tabular}{lcccl}
\toprule
\textbf{Domain} & \textbf{Ceiling} & \textbf{Models at ceiling} & \textbf{Total tested} & \textbf{Exceptions} \\
\midrule
Wikipedia & 1 & 6/6 & 6 & None \\
C4        & 2 & 8/9 & 9 & Mistral-7B ($\kstar\!=\!1$) \\
Code      & 2 & 4/5 & 5 & Pythia-1.4B ($\kstar\!=\!1$) \\
arXiv     & 3 & 5/5 & 5 & OLMo-1B ($\kstar\!=\!1$, capacity) \\
\bottomrule
\end{tabular}
\end{table}

\subsection{Cross-Model Transfer}
\label{sec:transfer}

To verify that $\kstar$ reflects properties of the generated text rather than scorer-specific artifacts, we conduct cross-scorer validation experiments.
Pythia$\leftrightarrow$OLMo cross-scoring retains 90--92\% classification accuracy; the 0v1 AUC exceeds 0.96 across all transfer pairs.
In a comprehensive cross-scorer sweep (12 generator$\times$scorer combinations on C4), 11 of 12 yield $\kstar\!=\!2$, with feature ordering Spearman $\rho > 0.90$.
These results confirm that depth features are text-intrinsic: they are imprinted by the generation process, not by the measurement apparatus.

% ===========================================================================
% 5  ANALYSIS
% ===========================================================================
\section{Analysis}

\subsection{The Domain Ceiling Framework}
\label{sec:domain_ceiling}

All observed $\kstar$ values can be explained by two factors:
(1) \emph{domain structural complexity}, which sets a ceiling on $\kstar$, and
(2) \emph{model domain-specific generation competence}, which determines whether a given model reaches the ceiling.

This framework generates falsifiable predictions.
For domains of intermediate complexity (C4, Code), we expect most models to reach the ceiling ($\kstar\!=\!2$) but allow for competence failures when pretraining data is mismatched.
For extreme domains, the framework predicts convergence: all domain-competent models should agree on $\kstar\!=\!1$ (Wikipedia) or $\kstar\!=\!3$ (arXiv), with model$\times$domain interaction disappearing.
Both predictions are confirmed empirically: Wikipedia shows zero exceptions (6/6), arXiv shows zero exceptions among domain-competent models (5/5), and C4/Code exceptions are systematically traceable to pretraining data composition (Mistral's web text bias, Pythia's code-sparse Pile, OLMo's limited academic coverage).

The JSD cliff position provides additional domain-specific information beyond $\kstar$.
Code and C4 share $\kstar\!=\!2$ but differ in where signal collapses: Code shows a 11--14$\times$ JSD cliff at the 1v2$\to$2v3 transition, while C4 shows a comparable cliff at the 0v1$\to$1v2 transition.
This indicates that domain structure shapes not only the depth of discriminability but also the pattern of information loss across generations.

\subsection{Representation Dependence}

Frozen neural probes (MLP on self-embeddings) yield model-dependent $\kstar$: Qwen achieves $\kstar_{\mathrm{neural}}\!=\!3$ (2v3 AUC = 0.6159, CI [0.610, 0.622]) while Pythia and OLMo remain at $\kstar_{\mathrm{neural}}\!=\!2$.
In contrast, our 15-dimensional model-free features produce consistent $\kstar$ across all models, validating the choice of model-free features for cross-model characterization.

\subsection{Generation Mode as a Confound}

Comparing continuation and rewrite (instruction-following) generation modes reveals a significant confound.
Rewrite mode amplifies the 1v2 AUC by +8.4 percentage points (0.696 vs.\ 0.611 on C4) through instruction-specific style artifacts, but cross-model transfer retention drops from 80--92\% (continuation) to 61--63\% (rewrite).
Continuation features form a universal subset of rewrite features (cont$\to$rw transfer: 89\% vs.\ rw$\to$cont: 82\%), supporting continuation mode as the canonical condition for generalizable cross-model claims.

\subsection{Naturalistic Validation on Common Crawl}
\label{sec:cc_validation}

To assess whether our controlled findings extend to real-world web data, we apply the depth classifier to 50{,}000 randomly sampled documents from Common Crawl 2025--2026 snapshots.
Table~\ref{tab:cc_depth} presents the predicted depth distribution.

\begin{table}[t]
\centering
\caption{Common Crawl predicted depth distribution with external correlates. Predicted depth correlates negatively with semantic coherence (AlignScore) and vocabulary richness (type-token ratio), consistent with recursive information loss.}
\label{tab:cc_depth}
\small
\begin{tabular}{lcccc}
\toprule
\textbf{Predicted depth} & \textbf{Proportion} & \textbf{Mean conf.} & \textbf{AlignScore} & \textbf{TTR} \\
\midrule
0 (human)  & 61.8\% & 0.724 & 0.682 & 0.571 \\
1          & 25.3\% & 0.481 & 0.641 & 0.534 \\
2          &  9.7\% & 0.387 & 0.618 & 0.512 \\
3+         &  3.2\% & 0.328 & 0.607 & 0.498 \\
\bottomrule
\end{tabular}
\end{table}

The distribution follows the expected heavy-tailed pattern: depth-0 dominates (61.8\%), depth-1 accounts for 25.3\%, and deeper generations are rare (12.9\% at depth $\geq 2$), consistent with the $\kstar\!=\!2$ prediction that beyond depth 2, information is saturated.
Classifier confidence drops sharply beyond depth 1 (0.724 $\to$ 0.481 $\to$ 0.387 $\to$ 0.328), reflecting the information-theoretic expectation that discrimination becomes progressively harder near and beyond $\kstar$.

\paragraph{External correlates.}
Table~\ref{tab:cc_correlations} shows that predicted depth correlates negatively with both AlignScore (semantic coherence; Spearman $\rho = -0.312$, $p < 10^{-50}$) and vocabulary richness measures (type-token ratio: $\rho = -0.274$; hapax ratio: $\rho = -0.238$).
These moderate correlations are consistent with recursive information loss reducing both semantic fidelity and lexical diversity, though the magnitudes reflect the expected noise in applying a controlled-setting classifier to heterogeneous web data.

\begin{table}[t]
\centering
\caption{Spearman correlation matrix for Common Crawl depth predictions and text quality indicators. All correlations significant at $p < 0.001$ ($n = 50{,}000$).}
\label{tab:cc_correlations}
\small
\begin{tabular}{lcccc}
\toprule
 & \textbf{Pred.\ depth} & \textbf{AlignScore} & \textbf{TTR} & \textbf{Hapax} \\
\midrule
Pred.\ depth & 1.000 & $-$0.312 & $-$0.274 & $-$0.238 \\
AlignScore   &       &  1.000   &  0.441   &  0.387   \\
TTR          &       &          &  1.000   &  0.812   \\
Hapax ratio  &       &          &          &  1.000   \\
\bottomrule
\end{tabular}
\end{table}

\paragraph{Domain-specific depth prevalence.}
Table~\ref{tab:cc_domains} breaks down depth predictions by Common Crawl content category.
Social media and user-generated content shows the highest synthetic proportion (56.2\% non-human), while technical documentation remains predominantly human-written (71.3\% depth-0).
A temporal comparison between January 2025 and January 2026 snapshots shows a 7.3 percentage point increase in predicted synthetic content (34.2\% $\to$ 41.5\%), consistent with the accelerating adoption of LLM-based content generation.

\begin{table}[t]
\centering
\caption{Common Crawl depth distribution by content category. Social/UGC content shows the highest synthetic contamination rate, while technical documentation remains predominantly human-written.}
\label{tab:cc_domains}
\small
\begin{tabular}{lcccc}
\toprule
\textbf{Category} & $N$ & \textbf{Depth-0} & \textbf{Depth-1} & \textbf{Depth-2+} \\
\midrule
News/journalism   & 12.4K & 52.1\% & 31.2\% & 16.7\% \\
Blogs/forums      & 15.8K & 58.4\% & 27.1\% & 14.5\% \\
Technical docs    &  8.2K & 71.3\% & 20.4\% &  8.3\% \\
Social/UGC        & 13.6K & 43.8\% & 34.7\% & 21.5\% \\
\bottomrule
\end{tabular}
\end{table}

\subsection{$\kstar$-Aware Data Curation}
\label{sec:curation}

We validate that domain-appropriate $\kstar$ predicts optimal curation granularity through a controlled curation experiment on C4 web text with synthetic contamination at 10\%, 30\%, and 50\% ratios.

For web text ($\kstar\!=\!2$), binary $\kstar$-aware filtering (retain depth 0--1, remove depth 2+) achieves 0.5941 MMLU accuracy, outperforming graduated depth weighting (best 0.5892 at $\alpha\!=\!0.9$), surprisal-based selection (0.5894), and DSIR distribution matching ($\approx$0.5867, near the no-filter baseline).
Dose-response analysis across three contamination ratios confirms that binary filtering maintains a consistent advantage, with HellaSwag improvements significant at all levels ($p \leq 0.004$), while graduated filtering collapses to no-filter baseline performance.

This result is predicted by the domain ceiling framework: since $\kstar\!=\!2$ for web text, there are only two actionable depth classes (depth $\leq 1$ and depth $\geq 2$), making binary filtering the natural granularity.
Graduated weighting introduces unnecessary complexity for a two-class problem, consuming model capacity on unreliable fine-grained depth estimates beyond the discrimination boundary.

\subsection{Threshold Sensitivity}

Our conclusions depend on the threshold $\theta = 0.60$.
A sensitivity analysis across $\theta \in [0.55, 0.70]$ reveals that the domain hierarchy is robust: arXiv maintains $\kstar\!=\!3$ and Wikipedia maintains $\kstar\!=\!1$ across the full range.
C4 $\kstar$ decreases from 2 (at $\theta\!=\!0.60$) toward 1 at higher thresholds ($\theta \geq 0.65$), reflecting the inherently weaker signal in homogeneous web text.
The sharp AUC discontinuity at boundary pairs (e.g., 1v2 AUC $\approx 0.62$ vs.\ 2v3 AUC $\approx 0.53$ for C4) means that $\kstar$ is stable across a wide range of reasonable $\theta$ values; only extreme choices would alter the qualitative conclusions.

% ===========================================================================
% 6  DISCUSSION
% ===========================================================================
\section{Discussion}

\paragraph{The one-wash laundering effect.}
Our central empirical finding is that recursive generation information loss is remarkably rapid.
For Wikipedia, a single generation pass suffices to erase all depth structure.
For web text and code, two passes reach the discrimination boundary.
Only richly structured academic text (arXiv) sustains discriminability to the third generation.
This ``one-wash laundering'' effect has direct implications for data governance: depth-aware curation beyond $\kstar$ offers no benefit, and the appropriate granularity is entirely determined by the source domain.

\paragraph{Why is $\kstar$ so small?}
The DPI guarantees that $\kstar$ is finite but provides no upper bound.
The empirical observation that $\kstar \leq 3$ across all tested conditions suggests that LLM generation is a strongly lossy channel: each pass erases a substantial fraction of the distributional features that distinguish one generation from the next.
Structured domains like arXiv sustain more generations because their distinctive formatting (LaTeX, citations, equations) provides additional signal dimensions that survive the generation process, but even this advantage extends $\kstar$ by only one step.

\paragraph{Implications for model collapse.}
The model collapse literature predicts progressive degradation under iterative self-retraining.
Our prompt-chain methodology measures a complementary quantity: the text-intrinsic information loss per generation pass without retraining effects.
The prompt-chain $\kstar\!=\!2$ for web text provides the operational boundary for data curation, as it reflects the minimum information loss achievable through recursive generation.
Retrain-chain experiments yield $\kstar\!=\!7$, but the additional five generations of discriminability arise from model-adaptation artifacts (vocabulary narrowing, repetition bias) that are specific to the retraining setup and do not generalize to wild contamination scenarios.

\paragraph{Limitations.}
Our model scale coverage extends to 32B but does not include 70B+ models, which remain important for future validation.
The curation benefit (C3) is validated only for C4 web text at specific contamination ratios; arXiv curation experiments yield null results, possibly due to evaluation metric insensitivity or insufficient training scale.
The Common Crawl validation lacks ground-truth depth labels and should be interpreted as a plausibility check rather than a precision estimate.
The scorer choice (Qwen2.5-1.5B) introduces a potential family bias, which we mitigate through cross-scorer validation but cannot fully eliminate.

% ===========================================================================
% 7  CONCLUSION
% ===========================================================================
\section{Conclusion}

We have presented the first systematic measurement of the generational discrimination boundary $\kstar$ for recursive synthetic text.
Grounded in the Data Processing Inequality, $\kstar$ captures the depth beyond which consecutive generations become statistically indistinguishable.
Through experiments spanning 9 model configurations (6 architecture families, 1B--32B), 4 source domains, and 4 decoding strategies, we establish that $\kstar \in \{1, 2, 3\}$, forming a domain hierarchy: Wikipedia $\kstar\!=\!1$, C4 and Code $\kstar\!=\!2$, arXiv $\kstar\!=\!3$.
Domain structural complexity is the sole systematic determinant of $\kstar$; model identity, scale, and decoding strategy modulate signal strength but not the domain ceiling.
The domain ceiling framework unifies all observed results, including apparent anomalies attributable to domain-specific competence failures.
Binary $\kstar$-aware filtering outperforms more granular alternatives for web text curation, confirming that $\kstar$ provides actionable guidance for data governance.
These findings establish the one-wash laundering effect: recursive generation erases depth structure within one to three passes depending on domain, with direct implications for the integrity of web-scale training data.

\bibliographystyle{plainnat}
\bibliography{references}

\end{document}
